\section{Parameters and arithmetic estimates}\label{sec:inputs}

Order the summands as $a\le b$. Write $H=\log c$ and $3x=\log R$.
All logarithms are natural. We fix $e=\min(\varepsilon,1/10)$ and
the parameters of Table~\ref{tab:parameters}. Every threshold for $x$
is uniform in the triple.
\fixedparameters

Every triple satisfies $H\le C_0R^{b_0'}$ for absolute $C_0>0$,
$b_0'\ge0$ \cite[Proposition~2.1]{Z29}. Taking exponent $2/9+1/18$
in \cite[Theorem~1.1]{Z29} gives $\log H\le 5x/6+O(1)\le x$
for large $x$.

Use the conventions of \cite[Section~2]{Z29}. A reduced fraction
$A/B$ with $B>0$ has height $h(A/B)=\log\max(|A|,B)$.
We normalize $v_p(p)=1$. Fix $\vartheta>0$. For $x\ge x_\vartheta$,
each signed quotient $z_0\ne\pm1$ of two sides and each prime $q\le R$ satisfy
\begin{equation}\label{eq:local}
 -\log|1-z_0|\le \exp(\vartheta x),\qquad
 v_q(1-z_0)\log q\le q\exp(\vartheta x).
\end{equation}
These bounds are \cite[Section~3, Fact~F2]{Z29}, from
\cite[Fact~F3, Lemma~2.1 and (2.8)]{ZhangThird}.

\begin{fact}\label{fact:bl}
Let $p$ be odd and let $\alpha_1,\alpha_2$ be multiplicatively
independent rational principal units at $p$. Put $\ell=\log p$ and
$T_i=\max(h(\alpha_i),\ell)$. For positive integers $b_1,b_2$, put
$b'=b_1/T_2+b_2/T_1$. Then
\begin{equation}\label{eq:bl}
 v_p(\alpha_1^{b_1}-\alpha_2^{b_2})\le
 \frac{24p T_1T_2}{(p-1)\ell^4}
 \max\{\log(b'\ell)+0.4,\,10\ell,\,10\}^2 .
\end{equation}
\end{fact}
This is the rational principal-unit case of
\cite[Corollary~2]{BL96}, used in \cite[proof of Proposition~6.1]{Z29}.

\begin{fact}\label{fact:primes}
Write $\vartheta(z)=\sum_{p\le z}\log p$. For $M,X\ge10^9$,
\[
 \vartheta(M)-\vartheta(M/2)>M/3,\qquad
 \vartheta(X)-\vartheta(10X/11)>X/12 .
\]
\end{fact}
Both inequalities follow by applying
$|\vartheta(z)-z|<0.2z/(\log z)^2$ at the two endpoints
\cite[Theorem~5.2]{Dusart}.
